\chapter{Explainable Artificial Intelligence (xAI)}

The term \emph{Explainable Artificial Intelligence} (xAI) refers to a set of processes and methods that allow human users to comprehend and trust the results and output created by machine learning algorithms. It is important for three main reasons:
\begin{enumerate}
    \item \ul{Trust}: Users need to understand why an AI makes a decision (e.g., credit denial or medical diagnosis), also depending on the peer group (social status).
    \item \ul{Regulation}: Laws such as the EU AI Act require traceability of AI systems.
    \item \ul{Debugging}: Identification of bias or errors in the model (e.g., discriminatory features).
\end{enumerate}

There are two main types of xAI methods:
\begin{itemize}
    \item \ul{Model-specific methods}: These methods are designed for specific types of models and leverage the internal structure of the model to provide explanations. Some examples include decision trees or attention maps for transformers.
    \item \ul{Model-agnostic methods}: These methods can be applied to any machine learning model and provide explanations based on the input-output behavior of the model, without relying on its internal structure.
\end{itemize}

Two important aspects of xAI are \emph{interpretability} and \emph{explainability}.
\begin{itemize}
    \item \ul{Interpretability}: The degree to which a human can understand the cause of a decision made by an AI model, which is often related to the complexity of the model and the transparency of its decision-making process.
    \item \ul{Explainability}: The degree to which an AI model can provide understandable explanations for its decisions.
\end{itemize}

Some important xAI methods include:
\begin{itemize}
    \item \ul{Grad-CAM (Gradient-weighted Class Activation Mapping)}: A technique that uses the gradients of a target concept, flowing into the final convolutional layer of a CNN, to produce a coarse localization map highlighting important regions in the image for predicting the concept.
    \item \ul{LIME (Local Interpretable Model-agnostic Explanations)}: A technique that explains the predictions of any classifier by approximating it locally with an interpretable model.
    \item \ul{SHAP (SHapley Additive exPlanations)}: A method that assigns each feature an importance value for a particular prediction, based on cooperative game theory.
    \item \ul{LRP (Layer-wise Relevance Propagation)}: A technique that decomposes the output of a neural network into contributions of its input features, providing insight into which features were most relevant for a given prediction.
\end{itemize}